% =====================================================================
%  Chapter 7 — Principle: Predictability (Leave Nothing to Chance)
%  Map: research/chapter-map.md (Chapter 7)
%  Shape: idea → origin → cost → what it opens up. Ideas, not how-to.
%  Rule: the author is never quoted in body text; short quotes only in
%  footnotes. Full statements live in interviews/.
% =====================================================================
\chaptersubtitle{Leave Nothing to Chance}
\chapter{Predictability}
\label{ch:memory}

Every program needs memory, and every language must decide who gives it back.
For fifty years there have been two main answers. In C, the programmer frees
memory by hand, and a single slip can corrupt a program or open it to attack.
In Java, Go, Erlang and most modern languages, a garbage collector frees it
automatically, and the program pays with pauses and overhead at moments no one
chose. Rust added a third answer, in which the compiler tracks ownership.
Silica's answer is different again, and it rests on a single demand: the
moment memory is released should be known in advance, every time.

\section{No one sweeping up behind}

Erlang, the language Silica learns most from, depends on a garbage collector.
When data might be shared, something must notice when the last user has gone,
and on Erlang's virtual machine that is the collector's job. Silica keeps
Erlang's independent actors and refuses the collector, so it must answer the
question another way.

The standard it set itself is strict. Releasing memory must be deterministic
and predictable, with nothing that behaves like a collector: no thresholds, no
sweeps, nothing done every so
often.\footnote{\src{silica/compiler/design_documents/data_structure_redesign_requirements.md}, line 66.}
When one of the project's AI agents once proposed a sweep that would run when
memory crossed a threshold, Barney refused it outright. The line is hard for a
reason that reaches beyond performance: if releasing memory is not
deterministic, it becomes impossible later to prove, formally, that the program
is correct (\S\ref{sec:proof}).

Silica's documents are equally unimpressed by the alternatives. Counting
references costs a little on every change. Rust's ownership rules they dismiss
as theatre,\footnote{\src{silica/docs/learn-silica.md}, line 27: ``No
borrow-checker theatre.''} and Barney, who knows Rust well from the bootstrap
compiler and the BEAM port, means the word almost literally. Getting every share
and transfer right in Rust feels like producing a musical in which the whole cast
is young and unstable. Given enough time, the show can go on.\footnote{Interview,
2026-11-13.} In Silica there is no cast to manage. Nothing is shared, everything
goes through the stack, and once data has been handed to a function, the
compiler will not let the old name be used again. Everything is predictable.

% TODO-INTERVIEW: Why write the bootstrap compiler in Rust anyway?

\section{Nothing outlives its home}

The first part of Silica's answer is that every piece of memory belongs
somewhere, and nothing that points into it may outlive the place it belongs.
The compiler refuses any program in which a reference could survive its home.
A dangling pointer---the root of a large share of the world's security
vulnerabilities---is not merely caught. It cannot be written.

\section{The heap nobody decided}

In October 2026 the project discovered that its runtime had not been doing what
its documents described. Since the first actors were built, data that had to
outlive a single step had been copied to a hidden store and released by
guesswork. The guesses were wrong in both directions: some memory was never
freed, and some was freed while still in use.\footnote{Commit
\code{19a7d19a48}, 27~October 2026.} Nobody had ever decided to build that
store.

What followed was two weeks of design at speed. Ownership by a single holder
was tried and withdrawn when measured. A store that never gave memory back was
adopted and abandoned: it was perfectly safe, and it would have grown by
something like 170 gigabytes an hour on a busy machine. Each rejected design
taught something about what the right one had to be.

\section{The heap and the file}

What settled the question was an insight about what the heap really is. The
memory a program borrows for long-lived data and the files it keeps on a disk
are usually treated as different worlds. Barney came to see them as two
expressions of the same thing, unsafe for the same reason: something else may
have changed the data since it was last stored. Beneath the details of how each
is implemented, using the heap and using persistent storage are the
same act.\footnote{Interview, 2026-11-13.}

Seen that way, the answer for memory was to stop keeping long-lived data in a
shared, unguarded place at all, and to keep each actor's data on that actor's
own stack. The answer for storage followed from the same insight. Silica's
local persistent storage uses its transport system
(Chapter~\ref{ch:transport}): data written to a disk and read back is checked and
converted exactly as data arriving from another machine would be. Safe transport
is safe transport, whether the data is persistent or not, local or remote.

% TODO-INTERVIEW: When did the heap/file insight come — during the October
%   fortnight, or earlier?

\section{A known end for everything}

Where it landed is simple to state. Each actor keeps everything it owns in its
own private space, and the compiler---not a background process, not a
heuristic---works out when each piece can be released. An idle actor holds no
working memory at all. Nothing is left to chance.

The new model was built across every desktop and server target by early
November 2026. It had a price: the compiler, itself written in Silica, briefly
grew too large to run on a Raspberry Pi, and had to be slimmed down again.

% TODO-INTERVIEW: Is "regions" still the right word for this principle now?

\section{Proof before code}
\label{sec:proof}

Determinism matters most for what it makes possible later. Proving a program
correct, in the way proof assistants such as Lean prove properties of code, is
a stated goal for Silica, in its specification and its development plan. A
proof is only possible if the program's behaviour is predictable, memory
included.

Barney imagines reversing the usual order. If the proof of correctness exists
before the code is written, and is expressed in a way language models can work
with, it becomes much easier for a model to write code that is correct. In that
arrangement the proof is the human's work. Models can help, but the proof is the
full responsibility of the person who is grounded in
reality.\footnote{Interview, 2026-11-13.}

% TODO-VERIFY: Where formal proof appears in the specification and the
%   development plan.

\section{What it costs}

The programmer must think about where data lives and how long it is needed,
questions a collector lets people ignore. Some structures must be copied where
other languages would share them. And the compiler does more work, which on
small machines is itself a resource to manage.

\section{What it opens up}

\begin{itemize}
  \item \textbf{Timing you can promise.} Without collector pauses, a program
        can respond within a known time---the foundation of control systems,
        audio, robotics and anything that must not stutter.
  \item \textbf{Tiny machines.} With no collector to feed and no hidden store
        to grow, a full actor system fits in half a megabyte.
  \item \textbf{Whole classes of vulnerability, gone.} Dangling pointers,
        double frees and use-after-free errors cannot be expressed, so they
        cannot be exploited.
  \item \textbf{Capacity you can plan.} Idle actors cost no working memory,
        so systems with enormous numbers of mostly quiet parts---sensors,
        connections, users---become affordable.
  \item \textbf{Programs that can be proved.} Deterministic memory keeps
        the door open to proving software correct before it is ever run.
\end{itemize}

None of these is the opportunity Barney cares about most. The one that matters
most is the one no one can yet imagine. What will curious, capable people build
on a foundation this stable? What systems will they make that are stable and
yet remarkably flexible? And which questions, too costly to answer before, will
they finally be able to answer?
